% !TEX root = ../main.tex
\chapter{Core ideas}\label{ch:core-ideas}
The rest of this manual leans on a small set of words. This chapter explains each one in
everyday terms, in the order you meet them when you use the app. \Cref{ch:glossary} has a
one-line definition of each, and of many more.

\section{Recordings, channels and spans}\label{sec:core-ideas-recordings}
A \term{recording}{recording} is one experiment's data file: the voltages from every
electrode in one fungus (or one set-up) over the whole time it was measured. A
\term{channel}{channel} is one electrode's trace inside it. The recording
\code{M2_aug_concat_fs1.mat}, for example, has 16 channels, each about 721 hours long, with one
value per second. A \term{span}{span} is a stretch of time on one channel: ``CH4\_A2, from
276.4 to 278.4 hours''.

\begin{plainterms}[A library of tapes]
Think of a recording as a long multi-track tape, one track per microphone. A channel is one
track. A span is a few minutes you cut out of one track to listen to closely.
\end{plainterms}

Every recording is stored with its sampling rate (how many values per second) and, where it
is known, its unit (volts or millivolts). Explore shows them on every page; the unit matters
whenever a setting is a threshold (\cref{ch:explore}). One recording is \term{held-out}{held
out}: it is kept unseen for the final test, and every page refuses it.

\section{Blocks, chains and templates}\label{sec:core-ideas-blocks}
A \term{block}{block} is one analysis step with one job: remove the slow drift, find
spikes, cut the signal into windows, group similar windows. The app has \ubnumblocks{}
blocks; \cref{ch:blocks} gives each a card. A block takes one kind of thing in and gives one
kind of thing out, and has a few settings.

A \term{chain}{chain} is blocks in a row, each feeding the next, starting from a span. You
build and run chains in Analyse. One position in a chain is a \term{step}{step}. Running a
chain over a span once is a \term{run}{run}; every run is numbered and kept.

A \term{template}{template} is a chain saved without its recording and span, so you can
apply the same steps somewhere else: on another span, or across many channels at once in
Discovery.

\begin{plainterms}[Kitchen words]
A block is one kitchen task: chop, boil, sieve. A chain is a recipe: those tasks in order. A
run is one time you cooked it, on one batch of ingredients. A template is the recipe card,
with no particular ingredients written on it.
\end{plainterms}

The app checks a chain as you build it. A block can only follow a block whose output it
accepts, so the \ui{Insert a stage} window greys out the blocks that do not fit at that point
and says why. The app also remembers each step's result: if you change the third block of a
four-block chain, the run button reads \ui{Re-run from 03}, and the first two steps' results
can be reused instead of recomputed.

\section{The seven kinds of thing a block passes on}\label{sec:core-ideas-types}
What passes between two blocks is always one of seven kinds of thing. These seven
\emph{types} are what decide which block may follow which. The project fixed the list at
seven on purpose: a method that fits none of them is out of scope, rather than a reason to
add an eighth (\code{Working/types/__init__.py}).

\begin{description}[style=nextline, leftmargin=1.6em, font=\sffamily\bfseries]
\item[Signal]
  A trace: one value per moment, plus how many moments per second. It is the raw span
  you start every chain with, and what a cleaning block gives back, cleaned. \emph{Made by}
  the source and by cleaning blocks such as \blockref{preprocessing.detrend} or
  \blockref{preprocessing.bandpass}. \emph{Read by} almost every block that comes first in a
  chain.
\item[Span set]
  A list of stretches of time, each with a start and an end, and optionally a label and a
  score: ``here, here and here''. It is what a detector's marks are, and what a
  person's marks are. \emph{Made by} detectors such as \blockref{detection.spike_v1} and
  \blockref{detection.threshold}. \emph{Read by} blocks that measure the marks, such as
  \blockref{interrogation.intervals}. A chain that ends in a span set writes its spans to the
  database as detections, ready for Review.
\item[Window set]
  The signal cut into pieces of a fixed length, like slicing a loaf, often with a table of
  measurements for each slice. \emph{Made by} \blockref{preprocessing.sliding_windows} and
  \blockref{preprocessing.window_matrix}. \emph{Read by} blocks that compare or group the
  slices, such as \blockref{catalogue.cluster}.
\item[Scores]
  One number for every moment of the signal: a running ``how unusual is this?'' line drawn
  under the trace. It lines up with the signal moment for moment, which is what lets a
  threshold turn it into a span set. \emph{Made by} \blockref{detection.matrix_profile}.
  \emph{Read by} \blockref{detection.threshold}.
\item[Encoding]
  The signal redrawn in another form: as a picture per window, or as a string of letters.
  \emph{Made by} encoders such as the symbolic encodings (\blockref{detection.sax_dsax}) and
  the Gramian images (\blockref{catalogue.gramian_gasf}). \emph{Read by}
  \blockref{detection.symbol_search} and \blockref{detection.drop_detection}, which search
  the letters.
\item[Grouping]
  A label for every window: which group it belongs to. \emph{Made by}
  \blockref{catalogue.cluster} (groups the machine finds) and
  \blockref{catalogue.manual_labels} (labels a person gave). \emph{Read by}
  \blockref{catalogue.classifier}.
\item[Model]
  A trained classifier: a program that has learned, from labelled windows, to label new
  ones. The type is only a pointer to the saved file. \emph{Made by}
  \blockref{catalogue.classifier}. Models are trained and compared in Models
  (\cref{ch:models}).
\end{description}

\begin{plainterms}[Why types at all]
Plugs and sockets. A kettle plug fits a kettle socket and not a phone charger. The seven
types are the plug shapes, and the app refuses to connect two blocks whose plugs do not fit,
before anything runs, rather than letting a long run fail an hour in.
\end{plainterms}

\begin{underthehood}[The types in code]
Each type is a small frozen dataclass in \code{Working/types/}: \code{Signal},
\code{SpanSet}, \code{WindowSet}, \code{Scores}, \code{Encoding}, \code{Grouping},
\code{Model}, each with \code{to_path} and \code{from_path} for the step cache. A
\code{Scores} is checked against the length of the \code{Signal} it came from; a
\code{WindowSet}'s feature table is deliberately not time-aligned. One chain-validation
function answers ``can this output feed this block'' for both the Insert window and the
runner (\code{Working/chain_validation.py}).
\end{underthehood}

\section{Detections and annotations: machine and human}\label{sec:core-ideas-detections}
Two kinds of mark can sit on a recording, and the app never lets them mix.
\begin{itemize}
  \item A \term{detection}{detection} is a span a \emph{machine} proposed: one mark written by
    one run of a chain. It is a \term{candidate}{candidate}, never a result.
  \item An \term{annotation}{annotation} is a span a \emph{person} marked, with that person's
    verdict. Explore's \ui{Take span for Review} writes one, and so does the project's older
    hand labelling.
\end{itemize}
They live in separate tables of the database. When you judge a detection in Review, your
verdict goes into a third table, of human verdicts on machine marks, and the detection itself
is untouched. No code path may write a person's verdict into a machine row, or the reverse:
this is rule 5 of the project (\code{CLAUDE.md}).

\begin{plainterms}[Why it matters]
One of the project's questions is ``where do people and the machine disagree?''
(\cref{sec:why-found}). That question has an answer only while the two sets of marks stay
apart. Mix them once, and you can no longer tell who said what: like a marked exam where the
teacher's ticks were written over the pupil's answers.
\end{plainterms}

\section{Verdicts}\label{sec:core-ideas-verdicts}
A \term{verdict}{verdict} is a person's judgement of one span. There are exactly five, the
same five for annotations and for judgements of detections, so there is no translation
table to drift (\code{Working/database/schema.py}, \code{VERDICTS}):
\begin{description}[style=nextline, leftmargin=1.6em, font=\sffamily\bfseries]
\item[seed] Worth keeping as an example to search for. In Review it also makes a new
  Library entry. Counts as accepted.
\item[interesting] Yes, this is a real pattern. Counts as accepted.
\item[not interesting] No. Counts as rejected.
\item[artifact] A fault of the recording, not the fungus (a loose wire, a shared ground).
  Kept, because one of the questions needs it, but kept out of training.
\item[unsure] You cannot tell.
\end{description}
In Review the keys are \texttt{S}, \texttt{I}, \texttt{N}, \texttt{A} and \texttt{U};
\texttt{Space} skips without writing anything. The list is fixed: the app never adds a sixth
verdict. Anything finer, such as a shape class (\ui{spike-train}, \ui{burst},
\ui{slow-drift}, \ui{plateau}), is a class or a tag given beside the verdict.

\section{Families, motifs, exemplars and seeds}\label{sec:core-ideas-families}
A \term{motif}{motif} is a shape that recurs: a particular dip, a sharkfin rise, a burst.
The Library is the catalogue of motifs the project has found, with every place each one
occurs (\code{docs/LIBRARY_STORAGE.md}, section~1). Each occurrence is a \emph{member}.

A \term{family}{family} is a group of members that look alike, made by grouping them by
shape. Two members of a family stand for it:
\begin{itemize}
  \item the \term{medoid}{medoid}, the member closest, on average, to all the others: the
    family's most typical member, picked by the computer;
  \item the \term{exemplar}{exemplar}, the member chosen to represent the shape.
    \unverified{A person can choose it, for example with the \ui{seed} verdict in Review;
    an entry imported from an older event store has its exemplar chosen by that import.}
\end{itemize}
A \term{seed}{seed} is a shape you search for: a Library exemplar, a family's medoid, or a
span you took in Explore. A \term{seed-search}{seed search} (in Discovery) slides the seed
along other channels and recordings and lists the closest look-alikes, which you then judge
in Review.

\begin{plainterms}[Families of faces]
Sort a box of photographs into piles of people who look alike. Each pile is a family. The
medoid is the face most like everyone else in the pile; the exemplar is the photo you would
pin on the wall to say ``this lot''. A seed search is holding one photo up against a crowd
and asking who looks most like it.
\end{plainterms}

\section{Nulls: is this more than chance?}\label{sec:core-ideas-nulls}
Any long, noisy trace contains look-alikes by accident. So a count on its own (``the chain
found 40 spikes'') says little. The question is always: \emph{more than chance would give?}

A \term{surrogate}{surrogate} is a fake copy of the signal that keeps its general character
(how jumpy it is, how much slow drift it has) but scrambles any real pattern in it. Run the
same chain on many surrogates and you learn what the chain finds in noise of the same kind.
That expectation is the \term{null}{null}. A result counts as more than chance only if the
real signal beats it.

\begin{plainterms}[Shuffling the deck]
You deal yourself a hand with three aces and wonder whether the deck is rigged. Shuffle and
deal a thousand more hands: if three aces turns up in a fifth of them, your hand proves
nothing. The surrogates are the reshuffled hands.
\end{plainterms}

In the app, \blockref{preprocessing.surrogate} makes surrogates. Analyse's \ui{surrogate
null} toggle (off unless you turn it on) runs a paired chain on one and shows detected
against surrogate counts. Discovery draws several surrogates per channel for its runs (20 for a template run, 200 for
a seed search, as the 3 October 2026 changes set them) and prints how many it drew next to
its counts. \Cref{app:nulls} explains each null the app
uses and what it does and does not test.

\begin{watchout}[A missing null is not a pass]
Some pages draw no null and say so in words. Read such a number as a description of what is
there, not as evidence that it beats chance.
\end{watchout}

\section{Jobs: work that takes a while}\label{sec:core-ideas-jobs}
Some work takes seconds; some takes hours. A \term{job}{job} is a piece of work the app runs
in the background while you carry on: a chain run, a Discovery run across channels, a model
training. The \ui{Jobs} button at the foot of the rail shows how many are running, and the
Jobs page lists what is in progress here, what is waiting on the cluster and what is waiting
in Review (\cref{ch:jobs}).

Work too long for your computer goes to the university's computing cluster. The app writes a
script for the cluster's queue (SLURM); the same analysis code runs there unattended, and
you bring its results back into the app (\cref{sec:why-built}).

\begin{plainterms}[Takeaway and dine-in]
A short run is a meal cooked at your table. A long one is an order sent to a big kitchen
down the road: you get a ticket (the job), carry on with your evening, and collect the order
when the ticket says it is ready.
\end{plainterms}
